\NormalFont\ShortTitle{2 Timóteo}
{\MT Segunda Carta de Paulo a Timóteo

\par }\ChapOne{1}{\PP \VerseOne{1}Paulo, apóstolo de Jesus Cristo, pela vontade de Deus, segundo a promessa da vida que está em Cristo Jesus,
\VS{2}a Timóteo,
{\ADD{meu}} amado filho. Graça, misericórdia,
{\ADD{e}} paz, da parte de Deus Pai, e de Cristo Jesus, nosso Senhor.
\VS{3}Agradeço a Deus, a quem desde os
{\ADD{meus}} antepassados sirvo com uma consciência pura, de que sem cessar tenho memória de ti nas minhas orações noite e dia.
\VS{4}Lembro-me das tuas lágrimas, e desejo muito te ver, para me encher de alegria;
\VS{5}Trago à memória a fé não fingida que há em ti,
{\ADD{fé}} que habitou primeiro na tua avó Loide, e na tua mãe Eunice; e tenho certeza que também em ti.
\VS{6}Por essa causa lembro-te de reacenderes o dom de Deus que está em ti pela imposição das minhas mãos.
\VS{7}Pois Deus não nos deu espírito de medo; mas sim de força, de amor, e de moderação.
\VS{8}Portanto não te envergonhes do testemunho do nosso Senhor, nem de mim, que sou prisioneiro dele; em vez disso, participa das aflições do Evangelho segundo o poder de Deus.
\VS{9}Ele nos salvou, e chamou com um chamado santo; não conforme as nossas obras, mas sim, conforme a sua própria intenção,
\FTNT{*}{{\FR 1:9 }
Ou: seu próprio propósito
} e a graça que nos foi dada em Cristo Jesus antes dos princípios dos tempos;
\VS{10}mas agora é manifesta pela aparição do nosso Salvador Jesus Cristo, que anulou a morte, e trouxe para a luz a vida e a imortalidade
\FTNT{†}{{\FR 1:10 }
Ou “incorrupção”, isto é, o apodrecimento do corpo
} por meio do evangelho;
\VS{11}do qual eu fui constituído pregador e apóstolo, e instrutor dos gentios.
\VS{12}Por causa disso também sofro essas coisas; porém não me envergonho. Pois sei em quem tenho crido, e tenho certeza de que ele é poderoso para guardar o que lhe confiei
\FTNT{‡}{{\FR 1:12 }
Ou: guardar o meu depósito
} até aquele dia.
\VS{13}Preserva o exemplo das sãs palavras que tens ouvido de mim, na fé e no amor que
{\ADD{há}} em Cristo Jesus.
\VS{14}Guarda o bem que te foi confiado
\FTNT{§}{{\FR 1:14 }
Lit. Guarda o bom depósito
} por meio do Espírito Santo que habita em nós.
\VS{15}Sabes isto, que todos os que estão na Ásia Romana me abandonaram, entre os quais estavam Fígelo e Hermógenes.
\VS{16}O Senhor conceda misericórdia à casa de Onesíforo, pois muitas vezes ele me consolou, e não se envergonhou de eu estar em cadeias;
\VS{17}Pelo contrário, quando veio a Roma, com empenho me procurou, e
{\ADD{me}} encontrou.
\VS{18}O Senhor lhe conceda que naquele dia ele encontre misericórdia diante do Senhor. E tu sabes muito bem como ele
{\ADD{me}} ajudou em Éfeso.

\par }\Chap{2}{\PP \VerseOne{1}Portanto tu, meu filho, fortifica-te na graça que há em Cristo Jesus.
\VS{2}E o que de mim ouviste entre muitas testemunhas, confia-o a pessoas fiéis, que sejam competentes para ensinar também a outros.
\VS{3}Tu, pois, sofre as aflições como bom soldado de Jesus Cristo.
\VS{4}Nenhum soldado em batalha se envolve com assuntos desta vida, pois tem como objetivo agradar àquele que o alistou.
\VS{5}E, também, se alguém está competindo como atleta, não recebe a coroa se não seguir as regras.
\VS{6}O lavrador que trabalha deve ser o primeiro a receber da partilha dos frutos.
\VS{7}Considera o que digo; que o Senhor te dê entendimento em tudo.
\VS{8}Lembra-te de Jesus Cristo, ressuscitado dos mortos, da descendência
\FTNT{*}{{\FR 2:8 }
Lit. “semente”
} de Davi, segundo o meu evangelho.
\VS{9}Por esse
{\ADD{evangelho}} sofro aflições, e até prisões, como um criminoso; mas a palavra de Deus não está presa.
\VS{10}Por isso tudo suporto por causa dos escolhidos, a fim de que também eles obtenham a salvação que está em Cristo Jesus com glória eterna.
\VS{11}Esta afirmação
{\ADD{é}} fiel: se morrermos com
{\ADD{ele}} , também com
{\ADD{ele}} viveremos;
\VS{12}se sofrermos, também com
{\ADD{ele}} reinaremos; se o negamos, também ele nos negará;
\VS{13}se formos infiéis, ele continua fiel; ele não pode negar a si mesmo.
\VS{14}Lembra-os disso, alertando-os diante do Senhor, que não tenham brigas de palavras, que não têm proveito algum,
{\ADD{a não ser}} para a ruína dos ouvintes.
\VS{15}Procura apresentar-te aprovado a Deus{\ADD{como}} um trabalhador que não tem de que se envergonhar, que usa bem a palavra da verdade.
\VS{16}Mas evita conversas profanas
{\ADD{e}} inútes; pois tendem a produzir maior irreverência.
\VS{17}E a palavra deles se espalhará como uma gangrena; entre os quais são Himeneu e Fileto.
\VS{18}Esses se desviaram da verdade, dizendo que a ressurreição já aconteceu; e perverteram a fé de alguns.
\VS{19}Porém o fundamento de Deus continua firme, e tem este selo:
\FTNT{†}{{\FR 2:19 }
Ou: esta inscrição
} o Senhor conhece os que são seus, e que todo aquele que faz uso do nome de Cristo, se afaste da injustiça.
\VS{20}Numa grande casa não há somente utensílios de ouro e de prata, mas também de madeira e de barro; e uns para honra, porém outros para desonra.
\VS{21}Portanto, se alguém se purificar destas coisas, será utensílio para honra, santificado e adequado para uso do Dono,
\FTNT{‡}{{\FR 2:21 }
Ou: Senhor
}
{\ADD{e}} preparado para toda boa obra.
\VS{22}Foge também dos desejos da juventude; e segue a justiça, a fé, o amor,
{\ADD{e}} a paz, com os que, de coração puro, invocam o Senhor.
\VS{23}E rejeita as questões tolas e sem instrução, como sabes que elas produzem brigas.
\VS{24}E ao servo do Senhor não convém brigar, mas sim ser manso com todos, apto para ensinar e suportar;
\VS{25}deve instruir com mansidão aos que se opõem, pois talvez Deus lhes dê arrependimento para conhecerem a verdade;
\VS{26}e se libertem da armadilha do diabo, em que foram presos à vontade dele.

\par }\Chap{3}{\PP \VerseOne{1}Sabe porém isto, que nos últimos dias virão tempos difíceis.
\VS{2}Pois haverá alguns que serão egoístas, gananciosos, presunçosos, soberbos, blasfemos, desobedientes a pais e a mães, ingratos, profanos,
\VS{3}sem afeto natural, irreconciliáveis, caluniadores, sem dominio próprio, cruéis, sem amor com os bons,
\VS{4}traidores, precipitados, orgulhosos, que amam mais os prazeres que a Deus.
\VS{5}Eles têm uma aparência de devoção divina, mas negam o poder dela. Afasta-te desses também.
\VS{6}Pois dentre esses são os que entram pelas casas, e levam cativas as mulheres insensatas carregadas de pecados, levadas com vários maus desejos;
\VS{7}que sempre estão aprendendo, e jamais conseguem chegar ao conhecimento da verdade.
\VS{8}E assim como Janes e Jambres se opuseram a Moisés, assim também esses se opõem a verdade; esses são corruptos de entendimento, e reprovados quanto à fé.
\VS{9}Eles, porém, não avançarão; pois a insensatez deles será evidente a todos, como também foi a daqueles.
\VS{10}Porém tu tens seguido a minha doutrina, conduta, intenção, fé, paciência, amor, perseverança,
\VS{11}perseguições, aflições; as quais me aconteceram em Antioquia, em Icônio e em Listra; tais perseguições sofri, e o Senhor me livrou de todas.
\VS{12}E também todos os que querem viver devotamente em Cristo Jesus sofrerão perseguição.
\VS{13}Mas os que são maus e enganadores irão de mal a pior, enganando e sendo enganados.
\VS{14}Tu, porém, continua nas coisas que aprendeste e das quais foste convencido, pois sabes de quem as aprendeste.
\VS{15}e que desde tua infância conheceste as Sagradas Escrituras, que podem te fazer sábio para a salvação pela fé em Cristo Jesus.
\VS{16}Toda a Escritura
{\ADD{é}} divinamente inspirada e
\FTNT{*}{{\FR 3:16 }
Ou: Toda a Escritura divinamente inspirada é também
} proveitosa para ensinar, para mostrar erros, para corrigir, e para instruir na justiça;
\VS{17}para que o homem de Deus seja completo, plenamente instruído para toda boa obra.

\par }\Chap{4}{\PP \VerseOne{1}Portanto te convoco diante de Deus e do Senhor Jesus Cristo, que irá julgar os vivos e os mortos em sua aparição e seu reino:
\VS{2}prega a palavra; insiste em tempo e fora de tempo; mostra os erros; repreende
{\ADD{e}} exorta, com toda paciência e ensino.
\VS{3}Pois virá tempo em que não suportarão a sã doutrina; em vez disso, terão comichão no ouvido e amontoarão para si instrutores conforme os seus próprios desejos maus;
\VS{4}desviarão os ouvidos da verdade, e se voltarão aos mitos.
\VS{5}Tu, porém, vigia-te em todas as coisas; suporta as aflições, faz a obra de evangelista, cumpre o teu serviço.
\VS{6}Pois já sou uma oferta de derramamento, e o tempo da minha partida está perto.
\VS{7}Lutei a boa luta, terminei a carreira, guardei a fé.
\VS{8}Desde agora a coroa da justiça me está reservada, a qual o Senhor, o justo Juiz, me dará naquele dia; e não somente a mim, mas também a todos os que amarem a sua vinda.
\VS{9}Procura vir logo até mim,
\VS{10}porque Demas me desamparou: ele amou o tempo presente, e partiu para Tessalônica; Crescente para a Galácia, e Tito para a Dalmácia.
\VS{11}Só Lucas está comigo. Toma Marcos, e o traz contigo, porque ele me é muito útil para o serviço.
\VS{12}Enviei Tíquico a Éfeso.
\VS{13}Quando vieres, traz a capa que deixei em Trôade na casa de Carpo, e os livros, principalmente os pergaminhos.
\VS{14}Alexandre, o que trabalha com cobre, me causou muitos males; o Senhor retribua a ele conforme suas obras.
\VS{15}Toma cuidado com ele também, porque ele se opôs muito às nossas palavras.
\VS{16}Na minha primeira defesa ninguém esteve comigo para me ajudar, pelo contrário, todos me desampararam. Que isso não seja levado em consideração contra eles.
\VS{17}Mas o Senhor me ajudou e me fortaleceu, a fim de que por mim a pregação fosse completamente divulgada, e todos os gentios
{\ADD{a}} ouvissem; e fui livrado da boca do leão.
\VS{18}E o Senhor me livrará de toda obra má, e me preservará para o seu reino celestial. A ele
{\ADD{seja}} glória para todo o sempre. Amém.
\VS{19}Cumprimenta Prisca e Áquila, e a casa de Onesíforo.
\VS{20}Erasto ficou em Corinto, e deixei Trófimo doente em Mileto.
\VS{21}Procura vir antes do inverno. Êubulo, Pudente, Lino, Cláudia e todos os irmãos te cumprimentam.
\VS{22}O Senhor Jesus Cristo seja com o teu espírito. A graça seja convosco. Amém.{\ADD{A segunda carta a Timóteo (o primeiro bispo escolhido para a igreja dos efésios) foi escrita em Roma, quando Paulo foi apresentado pela segunda vez ao imperador Nero}}
\par }